% part: intuitionistic-logic
% chapter: soundness-completeness
% section: decidability

\documentclass[../../../include/open-logic-section]{subfiles}

\begin{document}

\olfileid{int}{sc}{dec}

\olsection{નિર્ણેયતા}
નોંધો કે પૂર્ણતા પ્રમેયની સાબિતી દરેક
$\Gamma \Proves/ !A$ માટે અનંત સંખ્યામાં જગતો ધરાવતું
એવું નિદર્શ આપે છે જે $\Gamma \Entails/ !A$ની સાક્ષી
પૂરે છે. નીચેનું વિધાન બતાવે છે કે $\Entails !A$
સાબિત કરવા બધા સાન્ત નિદર્શો માટે
$\mSat{M}{!A}$ સાબિત કરવું પૂરતું છે (એટલે, સાન્ત
જગતગણ ધરાવતા નિદર્શો માટે).

\begin{thm}\ollabel{thm:decidability}
  જો $\Entails/ !A$ હોય, તો એવું સાન્ત નિદર્શ છે કે
  $\mSat/{M'}{!A}$.
\end{thm}

\begin{proof}
  ધારો કે $\mModel{M}=\tuple{W, R, V}$ એવું છે કે
  $\mSat/{M}{!A}$ અને $P$ એ~$!A$માં આવતા
  !!{propositional variable}sનો ગણ છે.
  $W' = \Setabs{[w]}{w \in W}$ લઈને, જ્યાં
  $[w] = \Setabs{p \in P}{w \in V(p)}$, $R'$ને
  ઉપગણ-સંબંધ અને
  $V'(p)=\Setabs{[w]}{p \in [w]}$ રાખીને
  $\mModel{M'}=\tuple{W', R', V'}$ વ્યાખ્યાયિત કરો.
  $W'$ સાન્ત ગણ છે અને $\mModel{M'}$ એ
  !!a{relational model} છે તે સ્પષ્ટ છે.

  $!A$ પર આગમન કરીને બતાવી શકાય છે કે
  \[
    \mSat{M}{!A}[w] \text{ ત્યારે અને માત્ર ત્યારે જ જ્યારે } \mSat{M'}{!A}[{[w]}]
  \]
  આ $P$માંથી જ !!{propositional variable}s ધરાવતા
  બધા !!{formula}s~$!A$ માટે સાચું છે. તેની સાબિતી
  વાચક માટે સ્વાધ્યાય તરીકે છોડી છે.
  \sourcecorrection{OLMD-096}{મૂળમાં આપેલી આ રચના
  નિષેધ માટે સત્યતા જાળવતી નથી. દાખલા તરીકે,
  $W=\{u,v\}$, $R$ ઓળખસંબંધ અને $V(p)=\{v\}$
  હોય, તો $u$ ખાતે $\lnot p$ સાચું છે. પરંતુ નવા
  નિદર્શમાં $[u]=\emptyset \subseteq \{p\}=[v]$,
  એટલે $[u]$થી $p$ સાચું હોય એવા જગત સુધી પહોંચી
  શકાય છે અને ત્યાં $\lnot p$ ખોટું છે. સાન્ત-નિદર્શ
  પ્રમેયની બીજી સાબિતી જરૂરી છે; આ દાવાની ખામી
  નોંધ્યા વિના તેને ચકાસેલી સાબિતી માનવી નહીં.}
\end{proof}

\begin{prob}
\olref[int][sc][dec]{thm:decidability}ની સાબિતી પૂરી
કરો: ફક્ત~$P$માંથી વિધાનાત્મક ચલો ધરાવતા બધા
!!{formula}s~$!A$ માટે $\mSat{M,w}{!A}$ ત્યારે
અને માત્ર ત્યારે જ જ્યારે $\mSat{M',[w]}{!A}$
તે બતાવો.
\end{prob}

\olref{thm:decidability} પરથી $\Entails !A$ છે કે નહીં
તે નક્કી કરવાની કલનવિધિ મળે છે.

\end{document}
